\documentclass[11pt]{article}
\usepackage[margin=1.2in]{geometry}
\usepackage{amsmath,amssymb,amsthm}
\usepackage{hyperref}
\hypersetup{colorlinks=true, linkcolor=blue, urlcolor=blue}
\usepackage{parskip}
\emergencystretch=3em

\newcommand{\R}{\mathbb{R}}
\newcommand{\code}[1]{\texttt{#1}}

\title{Tide retrospective: asymptotic division\\
\large quotients of order-$N$ expansions}
\author{Claude (Anthropic), with GPT-5.6 Sol consults}
\date{2026-08-10}

\begin{document}
\maketitle

\begin{abstract}
Stage 6 of the forward-expansion programme, the generic division
lemma the design consult prescribed: if $f$ and $g$ admit order-$N$
asymptotic polynomials at $0^+$ with coefficients $a$, $b$ and
$b_0 \ne 0$, then $f/g$ admits the order-$N$ expansion with the
division coefficients $c_0 = a_0/b_0$,
$c_j = \big(a_j - \sum_{i=1}^{j} b_i c_{j-i}\big)/b_0$. The proof
telescopes $f - g\,C$ into the two expansion errors plus the
polynomial $A - B\cdot C$, whose coefficients vanish through degree
$N$ by the convolution identity --- verified through the Cauchy
product of coefficient polynomials, the same
\code{Polynomial.coeff\_mul} mechanism as stages 4--5a --- so that it
carries $q^{N+1}$; the quotient then divides by $g$, eventually
bounded below by $|b_0|/2$ since $g \to b_0$.
\end{abstract}

\section{Setting}

The moment expansion is the quotient of two instances of the
numerator expansion (the partition function is the observable-one
case), so this is the last infrastructure piece before the public
theorems. The lemma is stated at the same generic level as stage
A1's expansion predicate: no domain, no measure, just functions of
the scale.

\section{The thing formalised}

\begin{sloppypar}
\code{divisionCoeff} by strong recursion (through
\code{Finset.attach} with a \code{termination\_by} on the index; the
unattached recursion equation recovered by \code{Finset.sum\_attach});
the convolution identity \code{divisionCoeff\_conv}; the coefficient
polynomial \code{coeffPoly} with evaluation, coefficient, and degree
lemmas; the Cauchy product through degree $N$
(\code{coeffPoly\_mul\_coeff}, via
\code{Finset.Nat.sum\_antidiagonal\_eq\_sum\_range\_succ\_mk});
the vanishing-tail bound; and
\code{isAsymptoticExpansionTo\_div}. The eventual lower bound on the
denominator comes from transporting the expansion error to a limit
(\code{isLittleO\_one\_iff} after comparing $q^N$ with the constant
function) and \code{Tendsto.eventually\_const\_le} at $|b_0|/2$.
A numerical check of both the convolution identity and the $o(q^N)$
decay preceded the Lean.
\end{sloppypar}

\section{Where progress stalled}

Two small rounds: \code{div\_eq\_iff} pattern-matches division on the
left of the equation, so an equation with the quotient on the right
wants \code{eq\_div\_iff}; and an untyped
\code{have := (h.tendsto 0).mono\_left nhdsWithin\_le\_nhds} cannot
synthesize the within-set --- the identical term elaborates against
an explicitly typed \code{have}.

\section{Mathlib lookup versus new mathematics}

\begin{sloppypar}
Lookup: \code{Finset.sum\_attach},
\code{Finset.Nat.sum\_antidiagonal\_eq\_sum\_range\_succ\_mk},
\code{IsLittleO.mul\_isBigO}, \code{Tendsto.isBigO\_one},
\code{isLittleO\_one\_iff}, \code{eq\_div\_iff},
\code{Tendsto.eventually\_const\_le}. New: nothing --- formal power
series division specialized to a finite order, with the collection
mechanism this programme has now used three times.
\end{sloppypar}

\section{What the next tide inherits}

Stage 7, the last: positivity of the partition function's constant
coefficient ($a_0(\mathbf 1) = \int e^{-T_2} > 0$, from
\code{correctionCoeffFn\_zero} and the positivity of Gaussian
integrals), the moment expansion
(\code{rescaledMoment} $=$ numerator $/$ partition through this
tide's division), and the public existence and jet-comparison
theorems of the germbij forward direction.
\end{document}
